% !TeX program = xelatex
% !TeX encoding = UTF-8
% ============================================================================
% Master working manuscript assembled from:
%   (i) the author-originated drafts main2-2.tex and chineseOrgans-2.tex,
%  (ii) the generated draft material supplied in the two pasted-text files,
% (iii) the paper-writing guide and example papers supplied by the author, and
%  (iv) the research structure, mathematical cautions, and work plan developed
%       in the accompanying discussion.
%
% IMPORTANT STATUS NOTE:
% This is a research-compilation document, not a submission-ready manuscript.
% Statements whose proofs or computations have not been independently checked
% are explicitly followed by % TODO: comments.  Generated draft material is
% retained only as research material and must be independently verified.
% ============================================================================

\documentclass[11pt, reqno]{amsart}

%==============================================================================
% Packages
%==============================================================================
\usepackage[utf8]{inputenc}
\usepackage{amsmath, amssymb, amsthm, mathtools}
\usepackage[margin=1.25in]{geometry}
\usepackage{enumitem}
\usepackage[dvipsnames]{xcolor}
\usepackage{hyperref}

% Additional packages needed by the compiled source material.
\usepackage[scheme=plain]{ctex}
\usepackage{graphicx}
\usepackage{booktabs}
\usepackage{array}
\usepackage{longtable}
\usepackage{microtype}
\usepackage{caption}
\usepackage{subcaption}
\usepackage{tikz}
\usetikzlibrary{arrows.meta,positioning,calc,decorations.pathreplacing}

% Hyperref setup for clean, clickable blue links (mimicking the Anthropic paper)
\hypersetup{
colorlinks=true,
linkcolor=NavyBlue,
citecolor=NavyBlue,
urlcolor=NavyBlue
}

\allowdisplaybreaks
\setcounter{tocdepth}{2}
\setcounter{secnumdepth}{3}

%==============================================================================
% Theorem Environments
%==============================================================================
\newtheorem{theorem}{Theorem}[section]
\newtheorem{lemma}[theorem]{Lemma}
\newtheorem{proposition}[theorem]{Proposition}
\newtheorem{corollary}[theorem]{Corollary}

\theoremstyle{definition}
\newtheorem{definition}[theorem]{Definition}
\newtheorem{example}[theorem]{Example}

\theoremstyle{remark}
\newtheorem{remark}[theorem]{Remark}

%==============================================================================
% Custom Macros
%==============================================================================
\newcommand{\R}{\mathbb{R}}
\newcommand{\N}{\mathbb{N}}
\newcommand{\Z}{\mathbb{Z}}
\newcommand{\Q}{\mathbb{Q}}
\newcommand{\C}{\mathbb{C}}
\newcommand{\F}{\mathbb{F}}
\newcommand{\eps}{\varepsilon}

\newcommand{\one}{\mathbf{1}}
\newcommand{\DM}{\mathrm{DM}}
\newcommand{\MB}{\mathrm{MB}}
\newcommand{\LB}{\mathrm{LB}}
\newcommand{\Ree}{\operatorname{Re}}
\newcommand{\Imm}{\operatorname{Im}}
\newcommand{\spec}{\operatorname{spec}}
\newcommand{\diag}{\operatorname{diag}}
\newcommand{\dd}{\,\mathrm{d}}
\DeclareMathOperator*{\argmax}{arg\,max}
\DeclareMathOperator*{\argmin}{arg\,min}

% Visible status marker for the working draft.  Formal submission versions
% should remove every use of this macro.
\newcommand{\statusnote}[1]{%
  \par\smallskip\noindent
  \textcolor{BrickRed}{\textbf{Working-draft status.} #1}%
  \par\smallskip
}

%==============================================================================
% Document Metadata
%==============================================================================
\title[Directed Cyclic Networks]{Synchrony, Wave Selection, and Robustness in Directed Cyclic Networks with Demand-Modulated Coupling\\[0.4em]
\textnormal{\large Master Working Manuscript and Research Blueprint}}

\author{Carter Feng}

%==============================================================================
% Document Body
%==============================================================================
\begin{document}

\begin{abstract}
Directed cyclic networks provide a tractable setting in which to study the
existence, stability, bifurcation, and robustness of synchronous states.  This
master document consolidates the author's original five-node models motivated
by the Wu-Xing diagram, the subsequent proposal to generalize the model to an
arbitrary number of nodes, the planned comparison with bilinear interaction
laws, and the proposed analyses of mode selection, Hopf bifurcation,
heterogeneity, and perturbation propagation.  Formulas and arguments already
recorded in the source drafts are typeset here, while all incomplete proofs,
derivations, computations, and literature checks are marked explicitly by
\texttt{TODO} comments in the source.  The intended final paper asks how
network size, interaction range, and nonlinear coupling law jointly determine
synchrony, wave selection, rotating waves, and sensitivity in directed cyclic
networks, with the five-node Wu-Xing-inspired network serving as the principal
application.
\end{abstract}

\maketitle

\keywords{cyclic networks, synchrony, circulant matrices, Fourier modes,
mode selection, Hopf bifurcation, rotating waves, heterogeneity, sensitivity,
Wu-Xing}

\subjclass[2020]{34C23, 34C15, 34D20, 37C80, 37G15, 92B25}

\tableofcontents
\newpage

%==============================================================================
\section*{Document Status and Provenance}
%==============================================================================

This document is intentionally broader than a submission-ready paper.  It is a
single formal repository for the mathematical model, the proposed theorem
package, the incomplete calculations, the numerical plan, the writing plan,
and the tangential ideas developed during the project.  The two files
\texttt{main2-2.tex} and \texttt{chineseOrgans-2.tex} are treated as the
author-originated research drafts.  The two pasted-text files contain generated
draft material and are retained only as a source of candidate organization,
formulae, and research questions.  No generated theorem or numerical claim is
to be treated as established until it has been independently rederived and
reproduced.

Throughout the manuscript, a \emph{synchronous state} is a mathematical state
in which all normalized node variables are equal.  Any interpretation of this
state as physiological balance is a modeling assumption rather than a clinical
claim.  The arbitrary-$n$ theory is a general network-dynamics problem; it is
not itself asserted to be a model of Traditional Chinese Medicine.

%==============================================================================
\section{Introduction}
\label{sec:introduction}
%==============================================================================

\subsection{Broad mathematical context}

Networks of interacting dynamical components may support equilibria in which
all components evolve identically, spatially patterned equilibria, traveling or
rotating waves, quasiperiodic motion, and chaos.  In a homogeneous cyclic
network, the cyclic symmetry organizes perturbations into discrete Fourier
modes.  The topology determines which spatial modes are available, while the
nonlinear interaction law determines their growth rates, the location of
stability boundaries, and the nonlinear states that emerge when synchrony is
lost.

The present project studies a family of directed cyclic ordinary differential
equations.  The original motivation is a five-node network with two directed
cycles: a nearest-neighbor generating cycle and a step-two controlling cycle.
The principal mathematical questions are whether trajectories remain in a
physically admissible region, when positive synchronous equilibria exist, which
Fourier mode loses stability first, what bifurcation occurs at the stability
boundary, and how the conclusions change when the nodes are not exactly
identical.

\subsection{The motivating five-node network}

The five-element or Wu-Xing diagram places five functional systems on a cycle
of generation, called Sheng (生), together with a pentagram of control, called
Ke (克).  In the original drafts, the five state variables were associated with
Wood/Liver, Fire/Heart, Earth/Spleen, Metal/Lung, and Water/Kidney.  The
resulting graph is mathematically a directed five-cycle together with a second
directed cycle of step length two.

The terminology from Traditional Chinese Medicine supplies a motivating
interpretation for several local interactions: nourishment from a mother to a
child, ordinary restraint from a controller to a target, over-acting, and
rebellious counter-action.  The mathematical paper, however, treats these as
idealized interaction rules between normalized functional components.  It does
not claim that the state variables are direct measurements of anatomical
organs, nor does it establish medical efficacy.

\subsection{Existing modeling approaches}

The source drafts distinguish discrete and continuous models.  Discrete models
represent each component by finitely many states, such as ON/OFF or
excess/normal/deficiency.  Continuous models assign each component a
real-valued activity and use ordinary differential equations.  Much of the
continuous literature described in the drafts uses bilinear, Lotka--Volterra,
or mass-action interaction terms of the form $x_i x_j$.

The project began from the observation that a bilinear positive input
$c x_i x_{i-1}$ vanishes when the receiving component $x_i$ is depleted.  The
author's alternative term
\[
    \frac{c x_{i-1}}{1+x_i}
\]
was intended to encode stronger demand when the receiving component is small.
For precision, the present document calls this a \emph{demand-modulated} or
\emph{receiver-modulated} input.  It is not a standard source-saturating
Holling type-II response, because it remains linear in the source variable
$x_{i-1}$.

Related work will ultimately be integrated into this Introduction and the
Discussion rather than placed in a separate ``Related Work'' section.  The
literature review must cover cyclic feedback systems, circulant matrices,
equivariant bifurcation, rotating waves, nearly identical networks,
heterogeneity, network Green functions, and the existing Wu-Xing models.

% TODO: Complete and verify the literature review.  For every cited Wu-Xing
% model, record the exact equations, indexing convention, parameter source,
% analytical results, numerical results, and limitations.

\subsection{Central research questions}

The intended paper is organized around the following questions.

\begin{enumerate}[label=\textup{Q\arabic*.},leftmargin=3em]
    \item For which parameter values does the demand-modulated cyclic network
    preserve nonnegativity and boundedness?

    \item When does a positive synchronous equilibrium exist, and how many
    synchronous equilibria can occur?

    \item For arbitrary network size $n$ and interaction step $q$, what is the
    exact Fourier spectrum of the linearization at synchrony?

    \item In the important step-two case $q=2$, which discrete wave number is
    selected as the least stable or first unstable mode?

    \item How do these conclusions compare with a topology-matched bilinear
    model and with the six-term bilinear model used in prior work?

    \item Which codimension-one bifurcations occur at the origin, at the fold
    of synchronous equilibria, and at a symmetry-breaking stability boundary?

    \item When a complex Fourier pair crosses the imaginary axis, is the Hopf
    bifurcation supercritical or subcritical, and does it produce stable
    rotating waves?

    \item How do small node-to-node parameter differences shift the
    equilibrium, stability boundary, Hopf threshold, and spatial response to a
    localized perturbation?
\end{enumerate}

\subsection{Proposed main message}

The proposed single message of the final paper is:

\begin{quote}
\emph{The cyclic architecture determines the available spatial modes, while
the nonlinear interaction law determines their stability; the same Fourier
structure organizes wave selection, Hopf-generated rotating waves, and the
response to heterogeneity and localized disturbances.}
\end{quote}

\subsection{Planned contribution package}

The strongest unified version of the paper combines five directions that were
originally considered separately:

\begin{enumerate}[label=(\roman*),leftmargin=2.5em]
    \item exact comparison of demand-modulated and bilinear interaction laws;
    \item complete steady-state and Hopf bifurcation analysis;
    \item persistence and robustness under small heterogeneity;
    \item an arbitrary-$n$ Fourier and mode-selection theorem; and
    \item equilibrium sensitivity and localized perturbation propagation.
\end{enumerate}

These contributions are intended to form one paper rather than five unrelated
extensions.  The arbitrary-$n$ theory provides the organizing framework, and
the five-node network supplies the detailed application.

\subsection{Organization of the master document}

Section~\ref{sec:preliminaries} introduces the normalized state variables, the
cyclic indexing, and the five-node interpretation.  Section~\ref{sec:model}
records the demand-modulated model, the matched bilinear model, the published
six-term comparator, and the heterogeneous extension.  The proposed theorem
package is organized in Sections~\ref{sec:global}--\ref{sec:heterogeneity}.
Section~\ref{sec:numerics} specifies the numerical and reproducibility plan.
Section~\ref{sec:planned-results} gives the intended results and figure order.
The appendices preserve the original motif models, the pathology extensions,
the notes on prior Lyapunov calculations, the claim audit, the complete work
plan, and the writing and publication strategy.

%==============================================================================
\section{Mathematical and Interpretive Preliminaries}
\label{sec:preliminaries}
%==============================================================================

\subsection{State variables and normalization}

For a network with $n$ nodes, let
\[
    x_i=x_i(t)\ge 0, \qquad i\in\Z_n,
\]
denote the dimensionless normalized activity of node $i$.  The normalization is
intended to place the components on a common reference scale.  Equality of the
normalized variables therefore means equal activity relative to the chosen
reference levels; it does not mean that distinct physiological systems possess
the same physical quantity.

The appearance of $1+x_i$ in the demand-modulated denominators presupposes that
$x_i$ has been nondimensionalized by a characteristic activity or demand
threshold.  A dimensional derivation should introduce an activity scale $X_0$
and a time scale $T_0$, set $x_i=X_i/X_0$, and derive the dimensionless
parameters.

% TODO: Write the complete dimensional model and nondimensionalization.  State
% the physical dimensions of every dimensional coefficient before scaling.

Let
\[
    \one=(1,\ldots,1)^{\mathsf T}\in\R^n.
\]
A state $x=x^*\one$ is called \emph{synchronous}.  In the five-node application,
a positive synchronous equilibrium is interpreted cautiously as a normalized
balanced state.

\subsection{Cyclic shifts and Fourier modes}

Let $P:\R^n\to\R^n$ denote the predecessor shift
\[
    (Px)_i=x_{i-1},
\]
with all indices interpreted modulo $n$.  The discrete Fourier modes are
\[
    v^{(m)}_j=e^{i\theta_m j},
    \qquad
    \theta_m=\frac{2\pi m}{n},
    \qquad
    m=0,\ldots,n-1.
\]
They satisfy
\[
    Pv^{(m)}=e^{-i\theta_m}v^{(m)}.
\]
The mode $m=0$ is the synchronous mode.  The nonzero modes are
symmetry-breaking modes.  Complex-conjugate modes occur in pairs
$m$ and $n-m$.

\subsection{The five-node index assignment}

For the motivating network, the indexing from the source drafts is
\begin{center}
\begin{tabular}{cll}
\toprule
Index & Element & Functional-system label from the source drafts \\
\midrule
$x_1$ & Wood (木)  & Liver (肝) \\
$x_2$ & Fire (火)  & Heart (心) \\
$x_3$ & Earth (土) & Spleen (脾) \\
$x_4$ & Metal (金) & Lung (肺) \\
$x_5$ & Water (水) & Kidney (肾) \\
\bottomrule
\end{tabular}
\end{center}

The generating cycle is
\[
    x_1\longrightarrow x_2\longrightarrow x_3\longrightarrow x_4
    \longrightarrow x_5\longrightarrow x_1,
\]
and the controlling cycle is
\[
    x_1\dashv x_3\dashv x_5\dashv x_2\dashv x_4\dashv x_1.
\]
Writing $\sigma(i)$ for the generating predecessor or mother,
$\kappa(i)$ for the control target, and $\rho(i)$ for the controller,
\begin{equation}
\label{eq:five-maps}
    \sigma(i)=i-1,
    \qquad
    \sigma^{-1}(i)=i+1,
    \qquad
    \kappa(i)=i+2,
    \qquad
    \rho(i)=i-2.
\end{equation}

\subsection{General step-$q$ cyclic architecture}

For arbitrary $n$, the five-node maps are generalized by
\[
    \sigma(i)=i-1,
    \qquad
    \kappa_q(i)=i+q,
    \qquad
    \rho_q(i)=i-q.
\]
The case $q=2$ recovers the original control step.  If $\gcd(n,q)=1$, the
step-$q$ edges form one directed cycle.  If $\gcd(n,q)>1$, they decompose into
several cycles.  The Fourier calculations remain meaningful in either case,
but the graph interpretation must record the decomposition.

\subsection{Classical relations retained as modeling motivation}

The original drafts list the following generating and controlling relations.
They are included here as modeling motivation; their precise classical sources,
translations, and scope require verification.

\paragraph{Generating relations.}
\begin{enumerate}[label=\arabic*.,leftmargin=2.5em]
    \item The mother nourishes the child, and the child draws upon the mother.
    \item Deficiency in the mother leads to deficiency in the child
    (母病及子).
    \item Excess in the mother leads to excess in the child.
    \item Excess activity in the child leads to overdraw from the mother
    (子盗母气).
    \item Sustained weakness in the child eventually weakens the mother.
\end{enumerate}

\paragraph{Controlling relations.}
\begin{enumerate}[label=\arabic*.,leftmargin=2.5em,start=6]
    \item Ordinary restraint: the controller suppresses its target.
    \item Over-acting (相乘, described in the drafts as 贼邪): an excessive
    controller or deficient target produces marked weakening of the target.
    \item Rebellious counter-acting (相侮, described as 微邪): an excessive
    target or deficient controller weakens the controller in the reverse
    direction.
\end{enumerate}

The original motif draft also records the treatment principles
\[
    \text{实则泻其子},
    \qquad
    \text{虚则补其母},
    \qquad
    \text{母能令子虚},
    \qquad
    \text{子能令母实}.
\]

% TODO: For every Chinese phrase, identify the exact primary or scholarly
% source, edition, page or chapter, Chinese wording, and a conservative English
% translation.  Distinguish classical statements from later interpretations.

\subsection{Core modeling assumptions}

The principal model begins from the following assumptions.

\begin{enumerate}[label=(A\arabic*),leftmargin=3em]
    \item Each state variable is a nonnegative normalized activity.
    \item In the homogeneous system, all nodes obey the same intrinsic law.
    \item The wiring is invariant under cyclic relabeling.
    \item The predecessor input increases with predecessor activity.
    \item The predecessor input is enhanced when the receiving node is
    depleted.
    \item Direct restraint requires both the restrained node and its controller
    to be active.
    \item Delays, stochastic forcing, intervention, and heterogeneity are
    omitted in the base model and introduced only in later extensions.
\end{enumerate}

%==============================================================================
\section{Model and Analytical Framework}
\label{sec:model}
%==============================================================================

\subsection{The original five-node demand-modulated model}

The coherent five-node model selected from the original drafts is
\begin{equation}
\label{eq:five-demand}
\dot{x}_i
=
\underbrace{a x_i-kx_i^2}_{\text{logistic self-regulation}}
+
\underbrace{\frac{c x_{\sigma(i)}}{1+x_i}}_{\text{demand-modulated predecessor input}}
-
\underbrace{\frac{g x_i}{1+x_{\kappa(i)}}}_{\text{target-dependent cost}}
-
\underbrace{h x_i x_{\rho(i)}}_{\text{controller-mediated restraint}},
\qquad i\in\Z_5,
\end{equation}
where $a,k,c,g,h>0$.

The baseline values appearing in the source drafts are
\begin{equation}
\label{eq:baseline-parameters}
    a=2,
    \qquad
    k=0.5,
    \qquad
    c=1,
    \qquad
    g=0.5,
    \qquad
    h=0.3.
\end{equation}

\begin{center}
\begin{tabular}{c p{9.2cm} c}
\toprule
Parameter & Working interpretation & Baseline \\
\midrule
$a$ & intrinsic self-production rate & $2.0$ \\
$k$ & self-limitation coefficient; uncoupled carrying capacity $a/k$ & $0.5$ \\
$c$ & strength of predecessor input & $1.0$ \\
$g$ & strength of the target-dependent cost & $0.5$ \\
$h$ & strength of direct controller-mediated restraint & $0.3$ \\
\bottomrule
\end{tabular}
\end{center}

\begin{remark}[Terminology]
Earlier generated drafts called both rational interaction terms ``saturating''
and identified them with Holling type-II responses.  The positive input in
\eqref{eq:five-demand} is not source-saturating, because it is linear in
$x_{\sigma(i)}$.  The present document therefore uses the more precise term
\emph{demand-modulated}.
\end{remark}

\subsection{The arbitrary-$n$, step-$q$ demand-modulated model}

The proposed generalization is
\begin{equation}
\label{eq:DMnq}
\boxed{
\dot x_i
=
 a x_i-kx_i^2
 +\frac{c x_{i-1}}{1+x_i}
 -\frac{g x_i}{1+x_{i+q}}
 -h x_i x_{i-q},
 \qquad i\in\Z_n.
}
\tag{$\DM_{n,q}$}
\end{equation}
The original five-node model is $\DM_{5,2}$.

The intended main theorem package treats general $n$ and $q$ where feasible,
then derives sharper mode-selection results for $q=2$.  The detailed numerical
application focuses on $n=5$ and $q=2$.

\subsection{A topology-matched bilinear comparator}

To isolate the effect of functional form without changing the graph or the
number of coefficients, the proposed matched bilinear analogue is
\begin{equation}
\label{eq:MBnq}
\boxed{
\dot y_i
=
 a y_i-k y_i^2
 +c y_i y_{i-1}
 -g y_i y_{i+q}
 -h y_i y_{i-q},
 \qquad i\in\Z_n.
}
\tag{$\MB_{n,q}$}
\end{equation}
This comparator is preferred for the main theoretical comparison because it
uses the same topology and the same number of parameters as
\eqref{eq:DMnq}.

\subsection{The six-term bilinear model from the prior literature}

The symmetric five-node model reproduced in the generated draft has the form
\begin{equation}
\label{eq:LB-five}
\dot y_i
=
 a y_i-b y_i^2
 +c y_i y_{i-1}
 -d y_i y_{i-2}
 -g_{\mathrm c} y_i y_{i+1}
 -h y_i y_{i+2},
 \qquad i\in\Z_5.
\end{equation}
Its proposed arbitrary-$n$, step-$q$ extension is
\begin{equation}
\label{eq:LBnq}
\dot y_i
=
 a y_i-b y_i^2
 +c y_i y_{i-1}
 -d y_i y_{i-q}
 -g_{\mathrm c} y_i y_{i+1}
 -h y_i y_{i+q},
 \qquad i\in\Z_n.
\tag{$\LB_{n,q}$}
\end{equation}
The six coefficients were described in the source drafts as gain,
saturation/self-limitation, is-generated, is-restricted, cost of generating the
child, and cost of restricting the target.

% TODO: Verify the exact equation, notation, and bibliographic metadata against
% the original Lin--Han publication rather than relying on the generated draft.

\subsection{Heterogeneous demand-modulated network}

The nodewise heterogeneous extension is
\begin{equation}
\label{eq:heterogeneous}
\dot x_i
=
 a_i x_i-k_i x_i^2
 +\frac{c_i x_{i-1}}{1+x_i}
 -\frac{g_i x_i}{1+x_{i+q}}
 -h_i x_i x_{i-q}.
\end{equation}
For perturbative analysis, write
\begin{equation}
\label{eq:hetero-expansion}
\begin{aligned}
    a_i&=a+\eps\alpha_i,
    &k_i&=k+\eps\kappa_i,
    &c_i&=c+\eps\chi_i,\\
    g_i&=g+\eps\gamma_i,
    &h_i&=h+\eps\eta_i.
\end{aligned}
\end{equation}
The paper must state clearly whether a coefficient is assigned to the receiving
node, the sending node, or the edge.  Equation~\eqref{eq:heterogeneous} adopts a
receiving-equation convention.

\subsection{Optional source-saturating variant}

A genuinely source-saturating and demand-modulated positive input would have
the form
\begin{equation}
\label{eq:true-saturation-candidate}
    c\,
    \frac{x_{i-1}}{K_{\mathrm s}+x_{i-1}}
    \frac{K_{\mathrm d}}{K_{\mathrm d}+x_i}.
\end{equation}
This expression is retained only as an optional future variant.  It is not part
of the frozen base model because adopting it would require redoing every
equilibrium, Jacobian, bifurcation, and sensitivity calculation.

% TODO: Decide permanently whether the submitted paper studies
% \eqref{eq:DMnq} or the true source-saturating variant.  Do not mix the two.

%==============================================================================
\section{Results I: Positivity, Boundedness, and Synchronous Equilibria}
\label{sec:global}
%==============================================================================

\subsection{Local existence and uniqueness}

\begin{proposition}[Planned local well-posedness statement]
For every initial condition in a neighborhood of the nonnegative orthant,
system~\eqref{eq:DMnq} has a unique local solution.  The vector field is smooth
on the region where $1+x_i>0$ and $1+x_{i+q}>0$ for every $i$.
\end{proposition}

% TODO: Supply the standard local-existence proof and state the maximal interval
% of existence precisely.

\subsection{Forward invariance and an absorbing box}

\begin{proposition}[Invariant nonnegative region and global bound]
\label{prop:absorbing-box}
For $a,k,c,g,h>0$, the nonnegative orthant is forward invariant under
\eqref{eq:DMnq}.  Moreover, the box
\[
    [0,X_{\max}]^n,
    \qquad
    X_{\max}=\frac{a+\sqrt{a^2+4kc}}{2k},
\]
is forward invariant and globally attracting.  Consequently, every
nonnegative solution exists for all $t\ge0$.
\end{proposition}

\paragraph{Existing proof outline from the drafts.}
At a boundary point with $x_i=0$, all terms in $\dot x_i$ vanish except the
nonnegative predecessor input, so
\[
    \dot x_i=cx_{i-1}\ge0.
\]
For
\[
    X(t)=\max_i x_i(t),
\]
a maximizing coordinate satisfies the comparison estimate
\[
    D^+X(t)\le aX-kX^2+c,
\]
where $D^+$ denotes the upper Dini derivative.  The right-hand side is negative
whenever $X>X_{\max}$.

% TODO: Write the complete proof using the upper Dini derivative, including the
% case of a nonunique maximizing index and the comparison argument establishing
% global existence and attraction of the box.

\subsection{The synchronous subspace}

The synchronous subspace
\[
    \mathcal S=\{x\in\R^n:x_1=\cdots=x_n\}
\]
is flow invariant for the homogeneous system.  Setting $x_i=x$ in
\eqref{eq:DMnq} gives
\begin{equation}
\label{eq:synchronous-scalar}
\dot x
=
 a x-(k+h)x^2+\frac{(c-g)x}{1+x}
=
\frac{x}{1+x}
\left[
(a+c-g)+(a-k-h)x-(k+h)x^2
\right].
\end{equation}

Define
\begin{equation}
\label{eq:A-mu}
    A=k+h,
    \qquad
    \mu=a+c-g.
\end{equation}
A positive synchronous equilibrium $x^*>0$ satisfies
\begin{equation}
\label{eq:synchronous-quadratic}
    A(x^*)^2+(A-a)x^*-\mu=0.
\end{equation}
The formal roots are
\begin{equation}
\label{eq:synchronous-roots}
    x_{\pm}
    =
    \frac{a-A\pm\sqrt{(a-A)^2+4A\mu}}{2A}.
\end{equation}

\subsection{Complete root classification}

The following classification records the complete case distinction identified
in the research plan.

\begin{proposition}[Candidate complete synchronous-equilibrium classification]
\label{prop:sync-classification}
Let $A=k+h>0$ and $\mu=a+c-g$.
\begin{enumerate}[label=(\roman*),leftmargin=2.5em]
    \item If $\mu>0$, equation~\eqref{eq:synchronous-quadratic} has exactly one
    positive root.

    \item If $\mu=0$, the roots are
    \[
        x=0,
        \qquad
        x=\frac{a-A}{A},
    \]
    and the second root is positive exactly when $a>A$.

    \item If $\mu<0$, there are two distinct positive roots exactly when
    \[
        a>A
    \]
    and
    \[
        a+c<g<
        a+c+\frac{(a-A)^2}{4A}.
    \]

    \item The two positive roots coalesce at
    \begin{equation}
    \label{eq:fold-location}
        g=g_{\mathrm f}
        :=a+c+\frac{(a-A)^2}{4A},
        \qquad
        x_{\mathrm f}=\frac{a-A}{2A}.
    \end{equation}
\end{enumerate}
\end{proposition}

% TODO: Supply a concise proof using the signs of the sum and product of the
% roots and the discriminant.  Check all boundary cases explicitly.

\subsection{The origin}

At $x=0$, the Jacobian of~\eqref{eq:DMnq} is
\begin{equation}
\label{eq:origin-jacobian}
    J(0)=(a-g)I+cP.
\end{equation}
Hence the Fourier eigenvalues are
\begin{equation}
\label{eq:origin-spectrum}
    \lambda_m(0)
    =(a-g)+c e^{-i\theta_m},
    \qquad m=0,\ldots,n-1.
\end{equation}
The spectral abscissa is $a+c-g=\mu$, so the origin is asymptotically stable
when $\mu<0$ and unstable when $\mu>0$.

% TODO: State and prove the corresponding proposition, including the
% nonhyperbolic case $\mu=0$.

\subsection{The matched bilinear synchronous state}

For~\eqref{eq:MBnq}, the synchronous scalar equation is
\[
    \dot y
    =a y-(k+g+h-c)y^2.
\]
Define
\begin{equation}
\label{eq:bilinear-M}
    M=k+g+h-c.
\end{equation}
When $M>0$, the positive synchronous equilibrium is
\begin{equation}
\label{eq:bilinear-sync}
    y^*=\frac{a}{M}.
\end{equation}
When $M<0$, the synchronous scalar equation contains a positive quadratic term
and admits finite-time blow-up for positive initial data.

% TODO: Prove a sufficient global boundedness condition for the full matched
% bilinear system, for example a condition of the form $k>c$, and distinguish it
% from the exact behavior on the synchronous subspace.

%==============================================================================
\section{Results II: Fourier Spectrum and Wave-Number Selection}
\label{sec:spectrum}
%==============================================================================

\subsection{Circulant linearization at synchrony}

Let $x^*>0$ be a synchronous equilibrium of~\eqref{eq:DMnq} and set
\begin{equation}
\label{eq:linear-coefficients}
    s=1+x^*,
    \qquad
    C_0=\frac{c}{s},
    \qquad
    G_0=\frac{g x^*}{s^2},
    \qquad
    H_0=h x^*.
\end{equation}
The diagonal entry is
\begin{equation}
\label{eq:d-diagonal}
    d
    =
    a-2kx^*
    -\frac{c x^*}{s^2}
    -\frac{g}{s}
    -h x^*.
\end{equation}
With the shift convention of Section~\ref{sec:preliminaries}, the Jacobian is
\begin{equation}
\label{eq:circulant-J}
    J
    =
    dI+C_0P+G_0P^{-q}-H_0P^q.
\end{equation}

\begin{theorem}[Candidate arbitrary-$n$, step-$q$ spectrum]
\label{thm:general-spectrum}
For $m=0,\ldots,n-1$, the Fourier eigenvalues of~\eqref{eq:circulant-J} are
\begin{equation}
\label{eq:general-eigenvalue}
\boxed{
    \lambda_m
    =
    d
    +C_0e^{-i\theta_m}
    +G_0e^{iq\theta_m}
    -H_0e^{-iq\theta_m}.
}
\end{equation}
Consequently,
\begin{align}
\label{eq:general-real}
    \Ree\lambda_m
    &=
    d+C_0\cos\theta_m+(G_0-H_0)\cos(q\theta_m),\\
\label{eq:general-imag}
    \Imm\lambda_m
    &=
    -C_0\sin\theta_m+(G_0+H_0)\sin(q\theta_m).
\end{align}
The synchronous equilibrium is locally asymptotically stable if and only if
\[
    \max_{0\le m<n}\Ree\lambda_m<0.
\]
\end{theorem}

% TODO: Supply the proof by diagonalizing the shift matrix with the unitary
% discrete Fourier transform.  Verify the sign convention against direct
% numerical diagonalization for random $n,q$ and parameter values.

\subsection{Normality of the homogeneous linearization}

Every circulant matrix is unitarily diagonalizable.  Therefore the homogeneous
Jacobian~\eqref{eq:circulant-J} is normal.  If all Fourier modes are stable,
then
\begin{equation}
\label{eq:normal-decay}
    \|e^{Jt}\|_2
    =
    e^{t\max_m\Ree\lambda_m}.
\end{equation}
Thus the homogeneous linearization has no nonmodal Euclidean-norm transient
growth.  Individual coordinates may still exhibit delayed or oscillatory
responses.

% TODO: Include a short proof and contrast this result with the nonnormal
% Jacobian that can arise after nodewise heterogeneity breaks circulancy.

\subsection{Quadratic dispersion relation for $q=2$}

For $q=2$, define
\begin{equation}
\label{eq:T-coefficient}
    T_0=G_0-H_0,
    \qquad
    u_m=\cos\theta_m.
\end{equation}
Using $\cos(2\theta)=2\cos^2\theta-1$, equation~\eqref{eq:general-real}
becomes
\begin{equation}
\label{eq:quadratic-dispersion}
    r_m:=\Ree\lambda_m
    =d-T_0+C_0u_m+2T_0u_m^2.
\end{equation}

\begin{theorem}[Candidate exact mode-selection theorem for $q=2$]
\label{thm:mode-selection}
Consider the discrete set
\[
    \mathcal U_n
    =
    \left\{\cos\frac{2\pi m}{n}:m=0,\ldots,n-1\right\}.
\]
\begin{enumerate}[label=(\roman*),leftmargin=2.5em]
    \item If $T_0\ge0$, the quadratic in
    \eqref{eq:quadratic-dispersion} is convex, and because $C_0>0$ its maximum
    over $\mathcal U_n$ occurs at $u=1$, corresponding to the synchronous mode
    $m=0$.

    \item If $T_0<0$, the quadratic is concave.  Its continuous maximizer is
    \begin{equation}
    \label{eq:u-star}
        u_*
        =-\frac{C_0}{4T_0}
        =\frac{C_0}{4(H_0-G_0)}.
    \end{equation}
    If $u_*\ge1$, the synchronous mode remains dominant.  If $0<u_*<1$, the
    dominant Fourier mode is any mode whose cosine value in $\mathcal U_n$ is
    closest to $u_*$.
\end{enumerate}
\end{theorem}

% TODO: Prove the theorem carefully for the discrete cosine set, handle ties,
% and characterize the parameter hypersurfaces on which two mode pairs have the
% same real part.

\subsection{Large-$n$ selected wave number}

Away from mode ties, the planned large-$n$ conclusion is
\begin{equation}
\label{eq:large-n-mode}
    \frac{m_*}{n}
    \longrightarrow
    \frac{1}{2\pi}\arccos u_*
    \qquad\text{as }n\to\infty.
\end{equation}
The anticipated eigenvalue error from replacing the continuous maximizer by the
nearest discrete wave number is of order $n^{-2}$ away from degeneracies.

% TODO: State and prove a precise corollary, including the assumptions needed
% to choose a unique representative from the conjugate pair.

\subsection{Five-node golden-ratio specialization}

For $n=5$ and $q=2$, let
\[
    \varphi=\frac{1+\sqrt5}{2}.
\]
The three distinct eigenvalue real parts are
\begin{align}
\label{eq:five-real-0}
    \Ree\lambda_0
    &=d+T_0+C_0,\\
\label{eq:five-real-1}
    \Ree\lambda_{1,4}
    &=d-\frac{\varphi}{2}T_0+\frac{1}{2\varphi}C_0,\\
\label{eq:five-real-2}
    \Ree\lambda_{2,3}
    &=d+\frac{1}{2\varphi}T_0-\frac{\varphi}{2}C_0.
\end{align}
The candidate threshold at which the pair $m=1,4$ overtakes the synchronous
mode is
\begin{equation}
\label{eq:five-threshold}
    H_0-G_0>\frac{C_0}{\varphi^2}.
\end{equation}

% TODO: Verify and prove the comparison among all three mode classes, not only
% the comparison between $m=0$ and $m=1,4$.

\subsection{Matched bilinear spectrum}

For the matched bilinear model~\eqref{eq:MBnq}, the positive synchronous state
$y^*=a/M$ has Fourier eigenvalues with real parts
\begin{equation}
\label{eq:matched-bilinear-real}
    \Ree\lambda_m^{\MB}
    =
    y^*
    \left[
        -k+c\cos\theta_m-(g+h)\cos(q\theta_m)
    \right].
\end{equation}
In particular,
\begin{equation}
\label{eq:matched-bilinear-uniform}
    \lambda_0^{\MB}=-a<0.
\end{equation}
Therefore, whenever the positive matched-bilinear synchronous equilibrium is
unstable, the instability is symmetry-breaking rather than synchronous.

For $q=2$,
\begin{equation}
\label{eq:matched-bilinear-quadratic}
    \frac{\Ree\lambda_m^{\MB}}{y^*}
    =
    -k+(g+h)+c u_m-2(g+h)u_m^2,
\end{equation}
whose continuous maximizer is
\begin{equation}
\label{eq:bilinear-u-star}
    u_{\MB,*}=\frac{c}{4(g+h)}.
\end{equation}

% TODO: Derive the full complex eigenvalue formula, prove the mode-selection
% rule, and identify all mode-tie boundaries.

\subsection{Six-term bilinear spectrum}

For the six-term model~\eqref{eq:LBnq}, the proposed real-part formula is
\begin{equation}
\label{eq:six-term-real}
    \Ree\lambda_m^{\LB}
    =
    y^*
    \left[
        -b+(c-g_{\mathrm c})\cos\theta_m-(d+h)\cos(q\theta_m)
    \right].
\end{equation}

% TODO: Verify the equilibrium and spectrum directly from the published model.
% Do not use this formula in the final paper until the original indexing has
% been checked.

%==============================================================================
\section{Results III: Comparison of Interaction Laws}
\label{sec:comparison}
%==============================================================================

\subsection{Structural comparison}

The principal controlled comparison is between~\eqref{eq:DMnq} and
\eqref{eq:MBnq}.  They have the same nodes, edge offsets, signs, and number of
parameters.  They differ only in the functional forms of the predecessor input
and target-dependent loss.

\begin{center}
\begin{tabular}{p{4.5cm}p{4.5cm}p{4.5cm}}
\toprule
Property & Demand-modulated model & Matched bilinear model \\
\midrule
Nonnegative orthant invariant & expected for all positive parameters & expected for all positive parameters \\
Global boundedness for every positive parameter choice & candidate theorem: yes & no; blow-up occurs on synchrony when $M<0$ \\
Synchronous-equilibrium equation & quadratic in $x^*$ & $y^*=a/M$ when $M>0$ \\
Fourier spectrum & explicit & explicit \\
Critical wave number for $q=2$ & candidate exact rule through $u_*$ & candidate exact rule through $u_{\MB,*}$ \\
Uniform mode at positive synchrony & branch dependent & always $-a$ \\
Symmetry-breaking Hopf possible & to be established & possible in selected parameter regimes \\
\bottomrule
\end{tabular}
\end{center}

\subsection{Functional-form interpretation}

For fixed predecessor activity, the bilinear input $c x_i x_{i-1}$ tends to zero
as the receiver $x_i$ tends to zero.  The demand-modulated input
$c x_{i-1}/(1+x_i)$ instead approaches $c x_{i-1}$.  This difference motivates
a figure comparing the two laws over the same normalized range.

% TODO: Decide on a fair normalization for the graphical comparison and state
% which quantities are held fixed.

\subsection{What should not be claimed}

The phrase ``bilinear synchrony is generically unstable'' appeared in generated
draft material but is not valid as a general statement.  The bilinear spectrum
has open stable and unstable regions.  The correct paper should provide exact
stability inequalities and should state instability only for specified
parameter sets or parameter regions.

Randomized stability fractions depend on the chosen prior over coefficients.
They may be included as a supplementary robustness experiment, but they should
not replace the exact theorem-based comparison and should not be presented as a
prior-independent property of either model.

%==============================================================================
\section{Results IV: Steady-State Bifurcations}
\label{sec:steady-bifurcations}
%==============================================================================

\subsection{Transcritical bifurcation at the origin}

At $\mu=0$, equation~\eqref{eq:synchronous-scalar} becomes
\[
    \dot x
    =
    \frac{x}{1+x}
    \left[(a-A)x-Ax^2\right].
\]
At the same parameter value, the nonsynchronous eigenvalues at the origin have
real parts
\[
    c(\cos\theta_m-1)<0,
    \qquad m\ne0.
\]
Thus only the synchronous mode is critical.

\begin{theorem}[Planned transcritical-bifurcation statement]
If $\mu=0$ and $a\ne A$, the origin and a synchronous equilibrium branch
undergo a generic transcritical bifurcation in the synchronous direction.
\end{theorem}

% TODO: Perform the center-manifold or scalar normal-form calculation and state
% the parameter orientation precisely.  Analyze the degenerate case $a=A$ as a
% codimension-two point.

\subsection{Saddle-node bifurcation of positive synchronous equilibria}

The fold candidate is given by~\eqref{eq:fold-location}.

\begin{theorem}[Planned saddle-node statement]
Assume $a>A$ and let $g=g_{\mathrm f}$.  If the synchronous zero eigenvalue is
simple and every nonsynchronous Fourier mode has nonzero real part, then the two
positive synchronous branches meet in a generic saddle-node bifurcation at
$x=x_{\mathrm f}$.
\end{theorem}

% TODO: Verify the full $n$-dimensional nondegeneracy conditions.  Do not infer
% full stability of the upper branch from the scalar synchronous equation alone.

\subsection{The actual bistable region}

The existence of two positive roots is not by itself sufficient for bistability
in the full system.  The actual bistable set is the subset of the two-root
region in which
\begin{enumerate}[label=(\roman*),leftmargin=2.5em]
    \item the origin is asymptotically stable;
    \item the upper positive synchronous equilibrium is asymptotically stable
    in every Fourier mode; and
    \item the lower branch has at least one unstable direction and separates
    the relevant basins.
\end{enumerate}

% TODO: Compute this subset analytically where possible and by continuation
% otherwise.  Determine whether the basin boundary is globally given by the
% lower branch or only locally organized by its stable manifold.

\subsection{Codimension-two steady-state degeneracy}

The point
\[
    \mu=0,
    \qquad
    a=A
\]
annuls both the linear and quadratic coefficients in the synchronous reduced
equation.  It is a candidate codimension-two degeneracy.

% TODO: Derive the local two-parameter unfolding and decide whether it belongs
% in the main paper or the supplement.

%==============================================================================
\section{Results V: Symmetry-Breaking Hopf Bifurcation and Rotating Waves}
\label{sec:hopf}
%==============================================================================

\subsection{Candidate Hopf boundary}

A symmetry-breaking Hopf candidate occurs when a nonzero mode $m_*$ satisfies
\begin{equation}
\label{eq:hopf-spectral}
    \lambda_{m_*}=i\omega_*,
    \qquad
    \lambda_{n-m_*}=-i\omega_*,
    \qquad
    \omega_*\ne0,
\end{equation}
while all remaining modes have negative real part.

The generated five-node draft identified a numerical crossing near
\[
    a=1.5,
    \qquad
    k=0.65,
    \qquad
    c=1,
    \qquad
    g=0.5,
    \qquad
    h\approx2.3718,
\]
with a critical pair reported as approximately
\[
    \pm0.310613i.
\]
This parameter point is retained only as a candidate starting point for an
independent continuation calculation.

% TODO: Reproduce this point from a clean implementation.  Confirm the
% equilibrium, all eigenvalues, the critical mode, and the crossing parameter.

\subsection{Hopf hypotheses}

For a rigorous local Hopf theorem, the final manuscript must verify:
\begin{enumerate}[label=(H\arabic*),leftmargin=3em]
    \item a unique simple conjugate pair lies on the imaginary axis;
    \item $\omega_*\ne0$;
    \item every other eigenvalue has nonzero real part;
    \item the crossing is transverse;
    \item the first Lyapunov coefficient is nonzero;
    \item no relevant low-order resonance invalidates the standard reduction.
\end{enumerate}

If $\nu$ is the chosen bifurcation parameter, transversality requires
\begin{equation}
\label{eq:transversality}
    \frac{\dd}{\dd\nu}\Ree\lambda_{m_*}(\nu_c)
    =
    \frac{\partial r_{m_*}}{\partial\nu}
    +
    \frac{\partial r_{m_*}}{\partial x^*}
    \frac{\dd x^*}{\dd\nu}
    \ne0.
\end{equation}

% TODO: Derive $\dd x^*/\dd\nu$ from the synchronous-equilibrium equation and
% obtain an explicit transversality formula for the selected parameter path.

\subsection{First Lyapunov coefficient}

After translating the equilibrium to the origin, write
\begin{equation}
\label{eq:Taylor-vector-field}
    F(z)
    =
    Jz
    +\frac12\mathcal B(z,z)
    +\frac16\mathcal C(z,z,z)
    +O(\|z\|^4).
\end{equation}
Let $v$ and $p$ satisfy
\[
    Jv=i\omega_*v,
    \qquad
    J^*p=-i\omega_*p,
    \qquad
    \langle p,v\rangle=1.
\]
The standard candidate formula recorded in the research plan is
\begin{align}
\label{eq:first-Lyapunov}
    l_1
    =
    \frac{1}{2\omega_*}
    \Ree\Big\langle p,
    &\mathcal C(v,v,\bar v)
    -2\mathcal B\bigl(v,J^{-1}\mathcal B(v,\bar v)\bigr)
    \\
    &+\mathcal B\bigl(\bar v,(2i\omega_*I-J)^{-1}
    \mathcal B(v,v)\bigr)
    \Big\rangle.
\end{align}

The cyclic Fourier structure is expected to reduce the calculation because
\[
    \mathcal B(v_m,\bar v_m)\text{ lies in mode }0,
    \qquad
    \mathcal B(v_m,v_m)\text{ lies in mode }2m,
\]
and
\[
    \mathcal C(v_m,v_m,\bar v_m)\text{ returns to mode }m.
\]
Thus only the modes $0$, $m_*$, and $2m_*$ should enter the scalar reduction.

% TODO: Compute the complete bilinear and trilinear tensors, reduce them in the
% Fourier basis, derive a reproducible expression for $l_1$, and determine its
% sign on a meaningful parameter interval.

\subsection{Supercritical and subcritical cases}

Under the standard orientation convention, $l_1<0$ corresponds to a
supercritical Hopf bifurcation with a small stable periodic branch, while
$l_1>0$ corresponds to a subcritical Hopf bifurcation with an unstable small
periodic branch.  A zero of $l_1$ is a candidate generalized Hopf or Bautin
point.

% TODO: State the sign convention used by the selected continuation software
% and reconcile it with the analytic formula.

\subsection{Rotating waves}

The leading-order form of a periodic branch associated with mode $m_*$ is
expected to be
\begin{equation}
\label{eq:rotating-wave}
    x_j(t)
    \approx
    x^*
    +R\cos\left(
        \omega_*t+\frac{2\pi m_*j}{n}+\phi
    \right).
\end{equation}
Adjacent nodes therefore differ by the time shift
\begin{equation}
\label{eq:phase-shift}
    \Delta t
    =
    \frac{m_*}{n}T,
    \qquad
    T=\frac{2\pi}{\omega_*}.
\end{equation}
If $\gcd(m_*,n)=1$, all $n$ nodes occupy distinct phases.  If
$\gcd(m_*,n)>1$, the wave has phase clusters.

% TODO: Apply the equivariant Hopf theorem with the exact cyclic group action,
% identify the isotropy subgroup, and verify the phase relation numerically on
% the continued periodic branch.

\subsection{Two-parameter continuation}

The planned continuation diagram should include the transcritical curve, the
fold curve, the Hopf curve, and mode-switch boundaries.  Candidate special
points include fold--Hopf intersections, Hopf--Hopf points, zero--Hopf points,
generalized Hopf points, and points at which two distinct Fourier modes tie.

% TODO: Continue the curves with AUTO-07p, MatCont, or COCO.  Analyze only the
% codimension-two points needed for the main story; place the remainder in the
% supplement or future work.

%==============================================================================
\section{Results VI: Heterogeneity, Sensitivity, and Perturbation Propagation}
\label{sec:heterogeneity}
%==============================================================================

\subsection{Persistence of a nearby equilibrium}

Let $x^{(0)}=x^*\one$ be a hyperbolic synchronous equilibrium of the homogeneous
system.  Under the perturbation~\eqref{eq:hetero-expansion}, the proposed
expansion is
\begin{equation}
\label{eq:heterogeneous-equilibrium-expansion}
    x(\eps)
    =
    x^*\one+\eps\xi+O(\eps^2).
\end{equation}

\begin{theorem}[Planned persistence theorem]
If $J=D_xF(x^*\one,0)$ is invertible, then for sufficiently small $|\eps|$
there exists a unique equilibrium near $x^*\one$, depending smoothly on
$\eps$.  If the homogeneous equilibrium is asymptotically stable, its stability
persists for sufficiently small heterogeneity.
\end{theorem}

% TODO: Supply the implicit-function-theorem proof and a precise stability
% persistence statement.

\subsection{First-order heterogeneous forcing}

At the homogeneous synchronous state, the first-order parameter forcing in the
receiving-equation convention is
\begin{equation}
\label{eq:heterogeneity-forcing}
    b_i
    =
    x^*\alpha_i
    -(x^*)^2\kappa_i
    +\frac{x^*}{1+x^*}\chi_i
    -\frac{x^*}{1+x^*}\gamma_i
    -(x^*)^2\eta_i.
\end{equation}
The first-order equilibrium correction satisfies
\begin{equation}
\label{eq:Jxi}
    J\xi=-b.
\end{equation}
Taking the discrete Fourier transform gives the candidate modal formula
\begin{equation}
\label{eq:xi-Fourier}
\boxed{
    \widehat\xi_m
    =
    -\frac{\widehat b_m}{\lambda_m}.
}
\end{equation}

% TODO: Derive the forcing carefully from the chosen heterogeneous convention
% and verify the $O(\eps^2)$ remainder numerically.

\subsection{Static sensitivity bound}

Since the homogeneous Jacobian is normal,
\begin{equation}
\label{eq:Jinverse-norm}
    \|J^{-1}\|_2
    =
    \frac{1}{\min_m|\lambda_m|}.
\end{equation}
The proposed first-order estimate is
\begin{equation}
\label{eq:static-sensitivity-bound}
    \|\xi\|_2
    \le
    \frac{\|b\|_2}{\min_m|\lambda_m|}.
\end{equation}
This estimate becomes large near a zero-eigenvalue bifurcation.  Near a Hopf
boundary, static sensitivity need not diverge because the critical eigenvalues
remain near $\pm i\omega_*$ rather than near zero.

% TODO: State the assumptions under which the norm equality and perturbation
% estimate hold, and distinguish static from frequency-dependent sensitivity.

\subsection{Localized equilibrium Green function}

If only the self-production coefficient at node $j_0$ is changed, then
\[
    b_j=x^*\delta_{j,j_0}.
\]
The proposed first-order equilibrium displacement is
\begin{equation}
\label{eq:green-function}
\boxed{
    \xi_j
    =
    -\frac{x^*}{n}
    \sum_{m=0}^{n-1}
    \frac{e^{i\theta_m(j-j_0)}}{\lambda_m}.
}
\end{equation}
Because the network is directed, the response need not satisfy
$\xi_{j_0+r}=\xi_{j_0-r}$.

% TODO: Derive the formula with the exact DFT normalization and define
% quantitative measures of directional response and spatial decay.

\subsection{Transient propagation from a localized perturbation}

For the homogeneous linearization and the initial condition
$y(0)=e_{j_0}$, the candidate response kernel is
\begin{equation}
\label{eq:transient-kernel}
\boxed{
    y_j(t)
    =
    \frac1n
    \sum_{m=0}^{n-1}
    e^{\lambda_m t}
    e^{i\theta_m(j-j_0)}.
}
\end{equation}
Near instability, the least stable Fourier mode should dominate the late-time
spatial pattern.

% TODO: Prove the formula, identify the propagation direction from phase
% velocities or peak times, and compare the linear prediction with the
% nonlinear system for several perturbation amplitudes.

\subsection{Harmonic response}

For a harmonic forcing of frequency $\Omega$, the modal transfer function is
recorded as
\begin{equation}
\label{eq:harmonic-response}
    \widehat y_m(\Omega)
    =
    \frac{\widehat f_m(\Omega)}{i\Omega-\lambda_m}.
\end{equation}
Near a Hopf boundary, forcing near
$\Omega=\Imm\lambda_{m_*}$ can produce a large dynamic response even when the
static inverse $J^{-1}$ remains moderate.

% TODO: Develop a precise frequency-response proposition and decide whether it
% belongs in the main paper or the supplement.

\subsection{A sufficient stability radius}

Let
\begin{equation}
\label{eq:stability-margin}
    \alpha
    =
    -\max_m\Ree\lambda_m>0.
\end{equation}
For a perturbed Jacobian $J_\eps=J+E_\eps$, a sufficient condition for all
eigenvalues to remain in the open left half-plane is expected to have the form
\begin{equation}
\label{eq:stability-radius}
    \|E_\eps\|_2<\alpha.
\end{equation}
A useful theorem requires an explicit bound
\[
    \|E_\eps\|_2
    \le
    K|\eps|\|\eta\|+O(\eps^2)
\]
that includes both the direct parameter perturbation and the shift of the
equilibrium.

% TODO: Derive a computable constant $K$ and state the resulting sufficient
% heterogeneity tolerance.

\subsection{Sensitivity of the Hopf threshold}

For a simple Hopf crossing, the proposed first-order shift of the critical
parameter is
\begin{equation}
\label{eq:hopf-threshold-shift}
    \delta\nu_c
    =
    -
    \frac{
        \Ree\langle p,E_\eta v\rangle
    }{
        \frac{\dd}{\dd\nu}\Ree\lambda_{m_*}
    }.
\end{equation}
This formula is intended to identify which spatial patterns of heterogeneity
advance or delay the instability.

% TODO: Derive the eigenvalue perturbation formula with the adopted complex
% inner-product convention and validate it against continuation of several
% heterogeneous examples.

\subsection{Optional persistence of the periodic orbit}

If the Hopf-generated periodic orbit is hyperbolic modulo its phase direction,
it should persist under sufficiently small heterogeneity, although exact
rotating-wave symmetry will be lost.

% TODO: State a precise periodic-orbit persistence result or cite the relevant
% theorem.  Confirm the prediction numerically using Floquet multipliers.

%==============================================================================
\section{Numerical Methods and Reproducibility}
\label{sec:numerics}
%==============================================================================

\subsection{Computational objectives}

The numerical work should verify and illustrate the analytical results rather
than replace them.  The required tasks are:
\begin{enumerate}[label=(\roman*),leftmargin=2.5em]
    \item verify the analytic Jacobians by finite differences;
    \item compare direct eigendecomposition with the Fourier formula;
    \item continue synchronous equilibrium branches and periodic orbits;
    \item compute transversality and the first Lyapunov coefficient;
    \item compute Floquet multipliers of periodic branches;
    \item validate first-order heterogeneous asymptotics;
    \item visualize localized static and transient response; and
    \item regenerate every figure and table from scripts.
\end{enumerate}

\subsection{Provisional numerical methods}

The generated draft records adaptive Runge--Kutta integration and LSODA in
stiff regions, with provisional tolerances
\[
    \mathrm{rtol}=10^{-9},
    \qquad
    \mathrm{atol}=10^{-11}.
\]
These values are retained as starting points rather than final choices.
Continuation may be performed with AUTO-07p, MatCont, or COCO.  The final paper
must state the software versions, parameter files, event definitions, and all
continuation settings.

% TODO: Select the production software stack and record exact versions,
% tolerances, continuation step sizes, branch labels, and stopping criteria.

\subsection{Validation tests}

The implementation should contain automated tests for:
\begin{enumerate}[label=(V\arabic*),leftmargin=3em]
    \item cyclic equivariance of the homogeneous vector field;
    \item analytic versus finite-difference Jacobians;
    \item equilibrium residuals;
    \item direct eigenvalues versus DFT eigenvalues;
    \item conjugacy $\lambda_{n-m}=\overline{\lambda_m}$;
    \item forward pointing of the vector field on coordinate faces;
    \item the absorbing-box differential inequality;
    \item agreement of analytic and numerical first Lyapunov coefficients;
    \item agreement of continuation points with direct simulation;
    \item quadratic accuracy of the heterogeneous equilibrium expansion; and
    \item the predicted phase shift on a rotating-wave branch.
\end{enumerate}

\subsection{Parameter studies}

The proposed network sizes are
\[
    n=5,\ 7,\ 9,\ 15,\ 31,
\]
with $q=2$ as the principal family.  The numerical atlas should include:
\begin{enumerate}[label=(\arabic*),leftmargin=2.5em]
    \item a stable synchronous regime;
    \item the transcritical crossing;
    \item the two-positive-equilibrium region;
    \item the saddle-node;
    \item a synchronous-mode instability, if present;
    \item a symmetry-breaking Hopf point;
    \item a rotating-wave periodic branch;
    \item a mode-switching boundary;
    \item a stable heterogeneous equilibrium; and
    \item a heterogeneity-induced shift or loss of stability.
\end{enumerate}

\subsection{Repository structure}

A proposed reproducible repository layout is
\begin{verbatim}
paper/
    manuscript/
    bibliography/
src/
    models/
    jacobians/
    fourier/
    bifurcation/
    lyapunov/
    sensitivity/
    validation/
figures/
data/
tests/
supplement/
README.md
environment.yml
\end{verbatim}

\subsection{Reproducibility statement placeholder}

% TODO: Replace this comment with a final reproducibility statement specifying
% the public repository, license, environment file, scripts generating each
% figure and table, random seeds, and instructions for rerunning continuation.

%==============================================================================
\section{Planned Results Presentation and Figure Order}
\label{sec:planned-results}
%==============================================================================

This section records the intended order in which the completed results should
appear.  It is not yet a results narrative.

\subsection{Result 1: Global admissibility}

% TODO: Present the local well-posedness, forward invariance, global bound, and
% complete synchronous-equilibrium classification.  Include the equilibrium
% bifurcation diagram and clearly separate existence from stability.

\subsection{Result 2: Arbitrary-$n$ Fourier spectrum}

% TODO: Present the circulant Jacobian, the general dispersion relation, and the
% exact $q=2$ mode-selection theorem.  Explain uniform versus symmetry-breaking
% modes before displaying the detailed inequalities.

\subsection{Result 3: Controlled comparison of interaction laws}

% TODO: Compare the demand-modulated model with the topology-matched bilinear
% model.  Include a functional-form figure and theorem-based stability maps.
% Relegate prior-dependent Monte Carlo comparisons to the supplement.

\subsection{Result 4: Steady and Hopf bifurcations}

% TODO: Present the transcritical and fold results, then the rigorously verified
% Hopf point, the first Lyapunov coefficient, periodic branch, and Floquet
% multipliers.

\subsection{Result 5: Heterogeneity and propagation}

% TODO: Present the implicit-function theorem, the Fourier sensitivity formula,
% the localized Green function, the transient response kernel, and the shift of
% the Hopf threshold.

\subsection{Result 6: Five-node synthesis}

% TODO: Specialize the arbitrary-$n$ theory to $n=5$, recover the golden-ratio
% expressions, and show how the general mode-selection and robustness theory
% organizes the original five-node model.

\subsection{Main figure plan}

\begin{longtable}{p{1.1cm}p{4.5cm}p{8.5cm}}
\toprule
Figure & Working title & Intended mathematical message \\
\midrule
\endfirsthead
\toprule
Figure & Working title & Intended mathematical message \\
\midrule
\endhead
1 & Network family & Show the directed predecessor cycle, the step-$q$ control edges, and the specialization $n=5,q=2$. \\
2 & Interaction-law comparison & Contrast the bilinear input and the demand-modulated input as the receiving state varies. \\
3 & Synchronous-equilibrium geometry & Display the origin, positive branches, transcritical point, fold, and the region with two positive roots. \\
4 & Dispersion and mode selection & Plot $\Ree\lambda_m$ against $m/n$ for several $n$ and compare with the continuous preferred wave number. \\
5 & Coupling-law stability comparison & Show stability and selected modes for the demand-modulated and matched bilinear systems on matched parameter planes. \\
6 & Two-parameter bifurcation organization & Show transcritical, fold, Hopf, and mode-switch curves. \\
7 & Hopf and rotating wave & Combine periodic-branch amplitude, period, Floquet multipliers, node time series, and phase-lag verification. \\
8 & Heterogeneity and localized sensitivity & Show the static Green function, transient heat map, first-order prediction, and $O(\eps^2)$ error. \\
9 & Five-node synthesis & Summarize how the general theory specializes to the Wu-Xing-inspired network. \\
\bottomrule
\end{longtable}

%==============================================================================
\section{Discussion}
\label{sec:discussion}
%==============================================================================

\subsection{Architecture versus interaction law}

The proposed conceptual conclusion is that topology and nonlinear interaction
law play distinct roles.  The cyclic topology fixes the Fourier modes and the
spatio-temporal symmetries available after bifurcation.  The coefficients and
functional forms determine the dispersion relation, the critical wave number,
the stability margins, and the nonlinear branch selected at instability.

\subsection{Interpretation of synchronous and symmetry-breaking modes}

The synchronous mode changes all normalized components together.  A nonzero
Fourier mode differentiates the components and generates a spatial pattern.
Within the motivating interpretation, this distinction suggests a mathematical
separation between uniform change and differentiated imbalance.  It does not,
by itself, establish a clinical interpretation for any particular mode.

\subsection{Why the arbitrary-$n$ extension matters}

The five-node golden-ratio expressions are consequences of the fifth roots of
unity.  The arbitrary-$n$ formulation shows which conclusions are genuinely
special to five nodes and which follow from cyclic symmetry more generally.
The $q=2$ quadratic dispersion relation also provides a direct bridge from
finite networks to a large-$n$ preferred wave number.

\subsection{Why heterogeneity belongs in the same paper}

The homogeneous Fourier spectrum controls not only stability but also the
leading response to parameter mismatch.  The same denominators $\lambda_m$
that identify weakly stable modes appear in the equilibrium sensitivity, Green
function, transient kernel, and Hopf-threshold shift.  Heterogeneity is
therefore not an unrelated extension; it is the natural robustness theory of
the homogeneous spectral result.

\subsection{Modeling limitations}

The final Discussion must state at least the following limitations.

\begin{enumerate}[label=(\arabic*),leftmargin=2.5em]
    \item The state variables are normalized abstract activities rather than
    direct physiological measurements.

    \item Exact $\Z_n$ symmetry and identical parameters are idealizations.

    \item A stable synchronous equilibrium is only one possible mathematical
    representation of balance; a bounded periodic state may also be compatible
    with normal operation in some applications.

    \item Linear stability determines only local behavior and does not give the
    basin of attraction.

    \item The model parameters have not been estimated from validated
    biological data.

    \item The model omits delays, noise, external forcing, and adaptive
    regulation in its base form.

    \item No medical or treatment conclusion follows from the mathematical
    analysis alone.
\end{enumerate}

\subsection{Relation to the example papers}

The supplied modeling examples suggest several presentation principles.  A
model should be explained in words before equations; the architecture should
be visible before the parameter analysis; the numerical procedure should be
reproducible; figures should communicate mechanisms rather than merely display
trajectories; and the Discussion should state both what the model explains and
what it omits.  These principles will guide the final rewrite.

\subsection{Future directions outside the main paper}

The following topics are intentionally reserved for later work:
\begin{enumerate}[label=(\roman*),leftmargin=2.5em]
    \item a rigorous slow--fast justification of local network motifs;
    \item a positivity-preserving model of all named Wu-Xing pathologies;
    \item adaptive homeostatic set points;
    \item treatment or intervention as an optimal-control problem;
    \item stochastic forcing and noise-induced switching;
    \item delay-induced instabilities;
    \item parameter estimation from validated measurements; and
    \item a separate critical reanalysis of previously reported Lyapunov
    spectra.
\end{enumerate}

%==============================================================================
\section{Conclusion}
\label{sec:conclusion}
%==============================================================================

% TODO: Write the final conclusion only after all theorem proofs, continuation
% diagrams, and heterogeneity calculations are complete.  It should state the
% question, the exact mathematical results, the principal consequence, and one
% highest-priority next direction without introducing new material.

%==============================================================================
\section*{Acknowledgments}
%==============================================================================

This work was conducted in the program SURE-DS 2026 (Summer Undergraduate
Research in Data Science 2026), organized by the Center of Scientific Machine
Learning in Material Science and supported by AFOSR grant FA9550-23-1-0725 and
the College of Science, Engineering, and Technology of Texas Southern
University.

% TODO: Add acknowledgments for the research mentor, readers, computational
% resources, and any additional grant support.

%==============================================================================
\appendix
%==============================================================================

\section{Original Motif Models from the Author Draft}
\label{app:motifs}

This appendix preserves the one-node, two-node, and three-node models from
\texttt{main2-2.tex}.  They are archival research material and are not part of
the frozen main model.

\subsection{One-node logistic motif}

The single-node motif was modeled by
\begin{equation}
\label{eq:one-node-logistic}
    \frac{\dd x}{\dd t}
    =
    \alpha x\left(1-\frac{x}{M}\right),
    \qquad
    \alpha,M>0.
\end{equation}

% TODO: Add the phase line and decide whether this elementary motif should
% appear in the final paper or only in pedagogical notes.

\subsection{Two-node bilinear mother--child motif}

The first two-node model in the author draft is
\begin{equation}
\label{eq:two-node-bilinear}
\begin{aligned}
    \dot x_1
    &=
    \alpha_1x_1\left(1-\frac{x_1}{M_1}\right)
    -\beta_1x_1x_2,\\
    \dot x_2
    &=
    \alpha_2x_2\left(1-\frac{x_2}{M_2}\right)
    +g_2x_1x_2.
\end{aligned}
\end{equation}
The draft records the candidate positive equilibrium
\begin{equation}
\label{eq:two-node-equilibrium}
\begin{aligned}
    x_1^*
    &=
    \frac{
        \alpha_1\alpha_2M_1
        -M_1M_2\beta_1\alpha_2
    }{
        \alpha_1\alpha_2
        +\beta_1g_2M_1M_2
    },\\
    x_2^*
    &=
    \frac{
        \alpha_1\alpha_2M_2
        +M_1M_2g_2\alpha_1
    }{
        \alpha_1\alpha_2
        +\beta_1g_2M_1M_2
    }.
\end{aligned}
\end{equation}
The draft proposes the condition $\alpha_1>M_2\beta_1$ for positivity of
$x_1^*$ and local stability of the positive equilibrium.

% TODO: Re-derive the equilibrium and stability result from scratch.  The
% original draft used $a_i$ and $\alpha_i$ inconsistently and contained an
% incomplete determinant calculation.

\subsection{Two-node demand-modulated motif}

A second two-node model was recorded as
\begin{equation}
\label{eq:two-node-demand}
\begin{aligned}
    \dot x_1
    &=
    \alpha_1x_1\left(1-\frac{x_1}{M_1}\right)
    -\frac{\beta_1x_1}{1+x_2},\\
    \dot x_2
    &=
    \alpha_2x_2\left(1-\frac{x_2}{M_2}\right)
    +\frac{g_2x_1}{1+x_2}.
\end{aligned}
\end{equation}

% TODO: Analyze the nullclines, positive equilibria, invariant region, and local
% and global stability of this motif if it is retained in a future motif paper.

\subsection{Piecewise threshold motif}

The source draft also proposed
\begin{equation}
\label{eq:two-node-threshold}
\begin{aligned}
    \dot x_m
    &=
    \alpha_mx_m(1-x_m)
    -\beta x_m\max(1-x,0),\\
    \dot x
    &=
    \alpha x(1-x)
    +g x_m\max(1-x,0).
\end{aligned}
\end{equation}

% TODO: Treat this as a piecewise-smooth system.  Determine the invariant
% regions, equilibria in each sector, switching manifold, and continuity or
% differentiability issues.

\subsection{Three-node motif}

The three-node diagram consists of a central node $x$, its generating
predecessor $x_m$, and a third node $x_c$ connected through the control motif.
The source draft records the candidate system
\begin{equation}
\label{eq:three-node-source}
\begin{aligned}
    \dot x_m
    &=
    \alpha_mx_m(1-x_m)
    -\beta x_m\max(1-x,0)
    -k\max(x_c-rx_m,0),\\
    \dot x
    &=
    \alpha x(1-x)
    +g x_m\max(1-x,0)
    -\beta x\max(1-x_c,0),\\
    \dot x_c
    &=
    \alpha_cx_c(1-x_c)
    +\mathcal G(x,x_c)
    -k\max(rx_m-x_c,0).
\end{aligned}
\end{equation}
The symbol $\mathcal G(x,x_c)$ is used here because the third positive term was
not unambiguously parsed in the source draft.

% TODO: Recover the intended third positive term from the author, define the
% piecewise-smooth sectors, and analyze the resulting three-node system.

\subsection{Local motifs and separation of time scales}

The original motif draft makes the conceptual observation that treatment rules
involving two or three nodes implicitly assume that the local subsystem is
approximately isolated over the relevant time interval.  A rigorous future
formulation would separate local variables $y$ from the remainder $z$ through
\begin{equation}
\label{eq:slow-fast-motif}
    \dot y=f(y,z),
    \qquad
    \dot z=\eps g(y,z),
    \qquad
    0<\eps\ll1.
\end{equation}

% TODO: Formulate and prove a singular-perturbation theorem under which the
% local motif tracks a stable branch $y=y^*(z)$ of the frozen subsystem.

\section{Extended Wu-Xing Pathology Models}
\label{app:pathologies}

\subsection{Mother consumption and child overdraw}

The source material distinguishes two different mechanisms:
\begin{enumerate}[label=(\roman*),leftmargin=2.5em]
    \item a deficient child receives stronger nourishment through the factor
    $(1+x_i)^{-1}$; and
    \item an excessive child consumes more from its mother through a separate
    pathological term.
\end{enumerate}
A candidate mother-loss term recorded in the generated draft is
\begin{equation}
\label{eq:overdraw-term}
    -\frac{g_d x_{i+1}x_i}{1+x_i}.
\end{equation}

\subsection{Rebellious counter-action}

A candidate counter-acting term from the source drafts is
\begin{equation}
\label{eq:wu-term-original}
    -\frac{b x_{i+q}}{1+\eps_0x_i}.
\end{equation}
The resulting extended five-node equation was recorded as
\begin{equation}
\label{eq:extended-pathology}
\dot x_i
=
 a_i x_i-kx_i^2
 +\frac{c x_{i-1}}{1+x_i}
 -\frac{g_d x_{i+1}x_i}{1+x_i}
 -\frac{g x_i}{1+x_{i+q}}
 -h x_i x_{i-q}
 -\frac{b x_{i+q}}{1+\eps_0x_i}.
\end{equation}

\begin{remark}[Positivity obstruction]
At $x_i=0$, the last term in~\eqref{eq:extended-pathology} does not vanish.
Indeed,
\[
    \dot x_i
    =
    c x_{i-1}-b x_{i+q}
\]
up to the terms that contain $x_i$.  This quantity can be negative.  Therefore
the nonnegative orthant is not automatically forward invariant for the
extended model as written.
\end{remark}

A positivity-preserving candidate mentioned in the research discussion is
\begin{equation}
\label{eq:wu-positive-candidate}
    -\frac{b x_i x_{i+q}}{K+\eps_0x_i},
\end{equation}
or a bounded product of receiver and target saturation factors.

% TODO: Select and justify a positivity-preserving counter-action mechanism.
% Recompute every equilibrium and Jacobian for the selected extension.  Do not
% claim that Proposition~\ref{prop:absorbing-box} transfers unchanged.

\subsection{Adaptive set points}

The source drafts proposed a contextual set point
\begin{equation}
\label{eq:adaptive-setpoint}
    \theta_i
    =
    \theta_i^0
    +\alpha_i x_{i-1}
    -\beta_i x_{i-q},
\end{equation}
leading to a self-regulation term
\begin{equation}
\label{eq:adaptive-logistic}
    k_i x_i(\theta_i-x_i).
\end{equation}

% TODO: Analyze positivity of the set point, invariant regions, synchronous and
% heterogeneous equilibria, and whether the cyclic Fourier structure survives
% in the homogeneous case.

\section{Notes on the Prior Six-Term Bilinear Model}
\label{app:lin-han}

\subsection{Reproduction targets from the generated draft}

The generated draft reports the following values for the six-term bilinear
model with all six coefficients equal to one:
\begin{equation}
\label{eq:reported-bilinear-eigs}
    -1,
    \qquad
    -0.5393\pm0.3919i,
    \qquad
    0.2060\pm0.6340i.
\end{equation}
The positive real part of the last pair makes that particular synchronous
equilibrium unstable.  The generated draft also reports a periodic orbit of
period approximately $17.626$.

% TODO: Reproduce every number from a clean implementation and the original
% published equation before including these values in a submitted paper.

\subsection{A structural zero Lyapunov exponent}

For a smooth autonomous system $\dot x=F(x)$, differentiation along a trajectory
gives
\[
    \frac{\dd}{\dd t}F(x(t))
    =
    DF(x(t))F(x(t)).
\]
Thus the flow direction is a solution of the variational equation.  A bounded
non-equilibrium trajectory is therefore expected to have a zero Lyapunov
exponent in the flow direction, subject to the standard regularity assumptions
used in the definition of the spectrum.

The generated draft records two published spectra with no zero exponent:
\begin{equation}
\label{eq:reported-spectra}
\begin{aligned}
    &\{0.3885,-0.1719,-0.3473,-0.6833,-0.9731\},\\
    &\{0.6700,0.6606,-0.0761,-0.0794,-1.1751\}.
\end{aligned}
\end{equation}

% TODO: Verify the original paper's numerical method, integration interval,
% transient removal, and exact interpretation of the reported numbers.

\subsection{Candidate periodic-orbit correction}

The generated draft reports a recomputed spectrum
\begin{equation}
\label{eq:recomputed-cycle-spectrum}
    \{0.0011,-0.1232,-0.2721,-0.2728,-0.9998\},
\end{equation}
which is consistent with one zero exponent and four negative exponents.

% TODO: Independently reproduce the orbit, period, tangent integration, QR
% reorthonormalization, convergence in averaging time, and the exponent values.

\subsection{Candidate boundary-equilibrium correction}

For a restriction-only parameter setting, the generated draft records the
Jacobian spectrum
\begin{equation}
\label{eq:boundary-equilibrium-spectrum}
    \{0,0,-1.3038,-2.3426,-3.6464\}
\end{equation}
at the state approximately
\[
    (0,0,0,2.3426,1.3038).
\]

% TODO: Verify the exact equilibrium, the dimension of the zero eigenspace,
% and the correct nonlinear stability classification.  Two zero eigenvalues do
% not by themselves establish neutral stability of the nonlinear equilibrium.

\subsection{Conservative subsystem and first integrals}

For the parameter restriction recorded in the drafts, the system reduces to
\begin{equation}
\label{eq:volterra-five}
    \dot x_i=x_i(x_{i-1}-x_{i+1}),
    \qquad i\in\Z_5.
\end{equation}
Two conserved quantities already derived in the draft are
\begin{equation}
\label{eq:first-integrals}
    I_1=\sum_i x_i,
    \qquad
    I_2=\prod_i x_i.
\end{equation}
Indeed,
\[
    \frac{\dd}{\dd t}\sum_i x_i=0,
    \qquad
    \frac{\dd}{\dd t}\log\prod_i x_i
    =
    \sum_i(x_{i-1}-x_{i+1})=0.
\]

A further candidate integral mentioned in the research discussion is
\begin{equation}
\label{eq:third-integral-candidate}
    I_3=\sum_i x_i x_{i+2}.
\end{equation}

% TODO: Verify $I_3$, determine functional independence of the integrals,
% connect the system to the periodic Volterra or Kac--van Moerbeke lattice, and
% do not claim quasiperiodicity on a torus without a complete integrability or
% invariant-torus argument.

\section{Claim Audit and Mathematical Cautions}
\label{app:claim-audit}

This appendix records the corrections and cautions that must govern every
future rewrite.

\begin{longtable}{p{5.1cm}p{8.7cm}}
\toprule
Claim appearing in earlier material & Required correction or qualification \\
\midrule
\endfirsthead
\toprule
Claim appearing in earlier material & Required correction or qualification \\
\midrule
\endhead
``The bilinear synchronous equilibrium is generically unstable.'' & False as a general statement.  Derive and state the exact Fourier stability inequalities.  Instability may be asserted only for specified parameters or regions. \\
``If $a+c\le g$, no positive synchronous state exists.'' & Incomplete.  When $a>k+h$, a two-positive-root interval can occur even with $a+c<g$. \\
``The whole two-root interval is bistable.'' & Not established.  The upper branch must be stable in every Fourier mode, and the basin geometry must be checked. \\
``A complex pair crossing the imaginary axis proves Hopf.'' & Incomplete.  Simplicity, nonzero frequency, transversality, and $l_1\ne0$ are required. \\
``The extended pathology model remains positive.'' & False for the counter-action term $-b x_{i+q}/(1+\eps x_i)$, which may point out of the nonnegative orthant at $x_i=0$. \\
``The positive input is Holling type II.'' & Inaccurate.  It is receiver-modulated but remains linear in the source. \\
``Two first integrals and small finite-time Lyapunov exponents prove quasiperiodicity.'' & Not sufficient.  Further integrability and invariant-torus analysis is required. \\
``Equal activity means equal organ function.'' & Equality is only meaningful after explicit normalization by separate reference levels. \\
``A negative spectral abscissa proves global recovery.'' & It proves local asymptotic stability only.  Basin size and global attractivity require additional arguments. \\
``Random parameter percentages measure model quality objectively.'' & The percentages are conditional on a chosen prior, validity filter, and classification procedure. \\
``The golden ratio is a physiological feature.'' & No.  It arises geometrically from the fifth roots of unity. \\
``The general $n$ model is a direct TCM model.'' & No.  It is a mathematical generalization motivated by the five-node topology. \\
\bottomrule
\end{longtable}

\section{Complete Research Work Plan}
\label{app:work-plan}

The phases below are ordered by mathematical dependency.  A later phase should
not be treated as complete until the stated gate from the earlier phase has
been passed.

\subsection{Phase 0: Freeze the paper contract}

\begin{enumerate}[label=0.\arabic*,leftmargin=3em]
    \item Adopt the mathematical title and central question.
    \item Freeze~\eqref{eq:DMnq} as the base model.
    \item Freeze~\eqref{eq:MBnq} as the matched comparator.
    \item Treat~\eqref{eq:LBnq} as a secondary literature application.
    \item Define normalized activities and the meaning of the unit denominator.
    \item Write a one-page assumption list.
    \item Create a claim ledger recording the support and status of every claim.
\end{enumerate}

\paragraph{Gate.}
Do not change the nonlinear model after beginning the Hopf and heterogeneity
calculations unless a genuine mathematical error forces a reset.

\subsection{Phase 1: Complete the literature and novelty audit}

Read and tabulate work in five clusters:
\begin{enumerate}[label=1.\arabic*,leftmargin=3em]
    \item cyclic feedback systems, rings, circulant spectra, and dispersion;
    \item coupled-cell synchrony and bifurcation;
    \item equivariant Hopf theory and rotating waves;
    \item robustness, nearly identical oscillators, heterogeneity, Green
    functions, and pseudospectra; and
    \item all existing Wu-Xing discrete and continuous models.
\end{enumerate}
For each paper, record the network, coupling law, generality in $n$, spectrum,
mode selection, Hopf coefficient, heterogeneity analysis, data source, and main
limitation.

\paragraph{Gate.}
Do not claim novelty or priority until the exact model and all proposed theorem
statements have been compared with the literature.

\subsection{Phase 2: Nondimensionalization and model specification}

\begin{enumerate}[label=2.\arabic*,leftmargin=3em]
    \item Write the dimensional model.
    \item Introduce activity and time scales.
    \item Derive the dimensionless equations.
    \item Define all graph maps and the role of $\gcd(n,q)$.
    \item State the homogeneous, heterogeneous, matched-bilinear, and six-term
    models.
    \item Produce a term-by-term interpretation table.
\end{enumerate}

\paragraph{Gate.}
Every term must have a sign, an edge direction, a parameter definition, and an
explicit statement of whether it is an assumption or a derivation.

\subsection{Phase 3: Build the reproducibility infrastructure}

Implement independently the vector field, analytic Jacobian, Fourier spectrum,
synchronous equilibria, finite-difference Jacobian, derivative tensors, ODE
solver, continuation interface, first Lyapunov coefficient, and heterogeneity
expansion.  Build the unit tests listed in Section~\ref{sec:numerics}.

\paragraph{Gate.}
No parameter map or continuation diagram should be trusted before the analytic
and numerical low-level tests agree.

\subsection{Phase 4: Prove the global foundation}

\begin{enumerate}[label=4.\arabic*,leftmargin=3em]
    \item Prove local existence and uniqueness.
    \item Prove forward invariance.
    \item Prove global boundedness with the Dini derivative.
    \item Prove invariance of the synchronous subspace.
    \item Complete the full root classification.
    \item Identify the origin, transcritical, fold, and degenerate parameter
    sets.
\end{enumerate}

\paragraph{Gate.}
The equilibrium classification must be complete and internally consistent
before any branch-stability statement is written.

\subsection{Phase 5: Prove the arbitrary-$n$ spectral theory}

\begin{enumerate}[label=5.\arabic*,leftmargin=3em]
    \item Express the Jacobian as a polynomial in the shift matrix.
    \item Derive the arbitrary-$n,q$ Fourier spectrum.
    \item Prove normality.
    \item Derive the $q=2$ quadratic dispersion relation.
    \item Prove exact discrete mode selection and characterize ties.
    \item Derive the large-$n$ selected wave number.
    \item Specialize to $n=5$ and verify the golden-ratio formulae.
\end{enumerate}

\paragraph{Gate.}
The theorem must predict the critical mode obtained by direct eigendecomposition
at every tested parameter point.

\subsection{Phase 6: Complete the controlled bilinear comparison}

\begin{enumerate}[label=6.\arabic*,leftmargin=3em]
    \item Prove positivity and equilibrium existence for the matched model.
    \item Prove a sufficient boundedness condition and document blow-up
    examples.
    \item Derive the arbitrary-$n,q$ spectrum and $q=2$ mode selection.
    \item Add the six-term literature model as a verified corollary.
    \item Use theorem-based comparisons in the main text.
    \item Move prior-dependent Monte Carlo comparisons to the supplement.
\end{enumerate}

\subsection{Phase 7: Prove the steady-state bifurcations}

\begin{enumerate}[label=7.\arabic*,leftmargin=3em]
    \item Prove the generic transcritical bifurcation.
    \item Analyze the $a=A$, $\mu=0$ degeneracy.
    \item Prove the full saddle-node nondegeneracy conditions.
    \item Determine the stability of each branch in all Fourier modes.
    \item Compute the actual bistable set rather than equating bistability with
    the two-root interval.
\end{enumerate}

\subsection{Phase 8: Complete the rigorous Hopf analysis}

\begin{enumerate}[label=8.\arabic*,leftmargin=3em]
    \item Choose a clean one-parameter path.
    \item Locate the critical point at high precision.
    \item Verify spectral simplicity and absence of relevant resonances.
    \item Prove transversality.
    \item Compute $\mathcal B$ and $\mathcal C$ symbolically and validate them.
    \item Reduce $l_1$ using Fourier selection rules.
    \item Determine the sign of $l_1$ on at least one meaningful interval.
    \item Continue the periodic branch and compute Floquet multipliers.
    \item Verify the rotating-wave phase relation.
    \item Continue the Hopf and fold curves in two parameters.
\end{enumerate}

\paragraph{Gate.}
Do not use the phrase ``nondegenerate Hopf bifurcation'' in the abstract until
transversality and $l_1\ne0$ have both been established.

\subsection{Phase 9: Develop the heterogeneity and sensitivity theory}

\begin{enumerate}[label=9.\arabic*,leftmargin=3em]
    \item Freeze the coefficient-assignment convention.
    \item Prove equilibrium persistence.
    \item Derive the first-order forcing and Fourier response.
    \item Derive the localized Green function.
    \item Derive the transient and harmonic kernels.
    \item Establish norm bounds and a computable stability radius.
    \item Derive the Hopf-threshold shift.
    \item Validate the $O(\eps^2)$ remainder.
    \item Test localized, low-mode, high-mode, random, and critical-mode
    heterogeneity.
\end{enumerate}

\subsection{Phase 10: Construct the numerical atlas}

Use theory to select representative parameter sets for every principal regime.
For every figure, record the parameters, solver, tolerances, initial condition,
integration horizon, continuation settings, and validation checks.  Repeat key
computations with tighter tolerances and an independent method.

\subsection{Phase 11: Write in dependency order}

\begin{enumerate}[label=11.\arabic*,leftmargin=3em]
    \item Finalize figures and captions.
    \item Write the theorem and results sections.
    \item Write the model and numerical-method sections.
    \item Write the Discussion and limitations.
    \item Write the Introduction.
    \item Write the Conclusion.
    \item Write the Abstract.
    \item Choose the final title.
    \item Remove repetition and every unsupported claim.
\end{enumerate}

\section{Manuscript Construction and Writing Guide}
\label{app:writing-guide}

\subsection{Recommended writing order}

The guiding framework supplied with the project recommends the order
\[
    \text{Results}
    \longrightarrow
    \text{Model}
    \longrightarrow
    \text{Discussion}
    \longrightarrow
    \text{Introduction}
    \longrightarrow
    \text{Abstract}
    \longrightarrow
    \text{Title}.
\]
This order forces the paper to be built around completed evidence rather than
around promises made too early.

\subsection{Single-message discipline}

The bottom-line conclusion should be repeated in the abstract, Introduction,
opening of the Results, and Discussion.  Supporting details should generally be
stated once.  The paper should read as a guided journey through one idea rather
than as a list of loosely connected calculations.

\subsection{Figures first}

The main figures and captions should be designed before the Results prose.  A
reader should be able to infer the central argument from the figures and their
captions.  Each caption should state the parameters, plotted quantity,
mathematical conclusion, and connection to a theorem.

\subsection{Model presentation}

The model should be explained in words before its equations.  State assumptions
clearly, display only the equations needed for the main story, and place
lengthy derivative tensors and routine expansions in an appendix or supplement.

\subsection{Results and interpretation}

The Results should distinguish theorem, computation, and observation.  Use
phrases such as ``we prove,'' ``numerical continuation shows,'' and ``simulation
suggests'' only when the corresponding evidence is present.  Interpretation
belongs mainly in the Discussion.

\subsection{Language and signposting}

Use simple language.  Each paragraph should carry one main message.  Technical
sections should begin with a sentence explaining what they establish and why a
reader might wish to continue.  Passive voice should not be used merely to
avoid ``we.''

\subsection{Claims and citations}

Every historical, medical, biological, and priority claim requires a citation
placed adjacent to it.  Classical Chinese statements require precise source
information.  Mathematical claims require a proof, a standard citation, or a
clear label as a conjecture or computational observation.

\subsection{Abstract template}

The final abstract should contain:
\begin{enumerate}[label=(\arabic*),leftmargin=2.5em]
    \item one or two sentences of background;
    \item one sentence stating the research question;
    \item two to four sentences stating what was proved or computed; and
    \item one sentence stating the broad mathematical significance.
\end{enumerate}
It should contain only results that are complete at submission time.

\section{Publication Strategy and Scope}
\label{app:publication}

\subsection{Primary target}

The current reach target is the \emph{SIAM Journal on Applied Dynamical Systems}
(SIADS).  A competitive submission should offer more than the observation that
a Jacobian is circulant.  The proposed strongest combination is:
\begin{quote}
exact discrete mode selection, an explicit and rigorously classified Hopf
bifurcation, and a Fourier Green-function theory for heterogeneous response.
\end{quote}

\subsection{Main-paper content}

The main manuscript should include:
\begin{enumerate}[label=(\roman*),leftmargin=2.5em]
    \item the general $n,q$ model and normalization;
    \item global positivity and boundedness;
    \item complete synchronous-equilibrium classification;
    \item arbitrary-$n,q$ Fourier spectrum;
    \item exact $q=2$ mode selection and the large-$n$ limit;
    \item the topology-matched bilinear comparison;
    \item transcritical, saddle-node, and one rigorous Hopf theorem;
    \item rotating-wave interpretation and periodic-orbit continuation;
    \item small-heterogeneity persistence and explicit sensitivity;
    \item a robustness bound and localized response; and
    \item the detailed $n=5,q=2$ application.
\end{enumerate}

\subsection{Supplementary content}

The supplement may contain full tensor formulae, expanded $l_1$ expressions,
extra continuation plots, additional $n,q$ maps, the complete six-term model,
Monte Carlo studies, random-heterogeneity covariance experiments, animations,
and all reproducibility details.

\subsection{Topics reserved for separate papers}

The motif slow--fast theory, the full pathology extension, adaptive set points,
optimal control, and a critical reanalysis of prior Lyapunov spectra should not
be forced into the same submission.

\subsection{Provisional page allocation}

\begin{center}
\begin{tabular}{p{8.5cm}r}
\toprule
Component & Pages \\
\midrule
Abstract and Introduction & $3$--$4$ \\
General network and coupling models & $3$ \\
Global dynamics and equilibria & $3$--$4$ \\
Fourier spectrum and mode selection & $5$--$6$ \\
Bilinear comparison & $2$--$3$ \\
Steady and Hopf bifurcations & $7$--$8$ \\
Heterogeneity, sensitivity, and propagation & $5$--$6$ \\
Numerical continuation and five-node application & $4$--$5$ \\
Discussion and Conclusion & $2$ \\
References & $2$--$3$ \\
\bottomrule
\end{tabular}
\end{center}

\section{Bibliographic Verification Checklist}
\label{app:bibliography-checklist}

Before submission, the bibliography should be checked for:
\begin{enumerate}[label=(\arabic*),leftmargin=2.5em]
    \item exact author names, titles, journal names, volumes, pages, and years;
    \item direct consultation of every paper whose equation is reproduced;
    \item precise primary or scholarly sources for the Chinese terminology;
    \item appropriate references for circulant matrices and the DFT;
    \item standard references for transcritical, saddle-node, and Hopf
    bifurcation;
    \item equivariant Hopf and rotating-wave theory;
    \item nearly identical networks, heterogeneity, and eigenvalue sensitivity;
    \item numerical continuation and first-Lyapunov-coefficient conventions;
    \item reproducible scientific computing and validated numerics if used; and
    \item removal of any citation that does not support the adjacent claim.
\end{enumerate}

%==============================================================================
% Bibliography
%==============================================================================

\begin{thebibliography}{99}

\bibitem{Davis1979}
P.~J. Davis,
\emph{Circulant Matrices},
Wiley, New York, 1979.

\bibitem{GolubitskyStewartSchaeffer1988}
M. Golubitsky, I. Stewart, and D.~G. Schaeffer,
\emph{Singularities and Groups in Bifurcation Theory, Vol. II},
Applied Mathematical Sciences, Vol.~69, Springer, 1988.

\bibitem{GolubitskyStewart2002}
M. Golubitsky and I. Stewart,
\emph{The Symmetry Perspective: From Equilibrium to Chaos in Phase Space and
Physical Space},
Birkh\"auser, 2002.

\bibitem{Huguet2014}
G. Huguet, J. Rinzel, and J.-M. Hup\'e,
Noise and adaptation in multistable perception: Noise drives when to switch,
adaptation determines percept choice,
\emph{Journal of Vision} \textbf{14} (2014), no.~3, Article~19.

\bibitem{WangKilpatrickJosic2020}
Y. Wang, Z.~P. Kilpatrick, and K. Josi\'c,
A hierarchical model of perceptual multistability involving interocular
grouping,
working manuscript, 2020.

\bibitem{Milo2002}
R. Milo, S. Shen-Orr, S. Itzkovitz, N. Kashtan, D. Chklovskii, and U. Alon,
Network motifs: Simple building blocks of complex networks,
\emph{Science} \textbf{298} (2002), 824--827.

\bibitem{Alon2007}
U. Alon,
Network motifs: Theory and experimental approaches,
\emph{Nature Reviews Genetics} \textbf{8} (2007), 450--461.

\bibitem{Kaptchuk2000}
T.~J. Kaptchuk,
\emph{The Web That Has No Weaver: Understanding Chinese Medicine},
revised edition, McGraw--Hill, 2000.

\bibitem{Porkert1990}
M. Porkert,
\emph{The Theoretical Foundations of Chinese Medicine},
MIT Press, 1990.

\bibitem{Lin2020}
% TODO: Replace this placeholder with the full verified bibliographic entry.
Lin and Han,
\emph{[Title and publication details to be verified]},
2020.

\bibitem{Nian2025}
% TODO: Replace this placeholder with the full verified bibliographic entry.
Nian and Li,
\emph{[Title and publication details to be verified]},
2025.

\bibitem{NeiJing}
% TODO: Supply a precise edition, translator, chapter, and page references.
\emph{Huangdi Neijing} (黄帝内经),
classical source; edition and translation to be specified.

\bibitem{NanJing}
% TODO: Supply a precise edition, translator, chapter, and page references.
\emph{Nanjing} (难经),
classical source; edition and translation to be specified.

\end{thebibliography}

\end{document}
